% I. Background and context (problem area and motivation) 
% II. However (problem description)
% III. So what? (problem dimension) 
% IV. Deficiencies with previous approaches
% IV. Solution

\section{Introduction} \label{sec:INTRODUCTION}

Many tools and development environments try to shorten development cycles to provide feedback as early as possible.
The \ac{lk}\footnote{\url{http://lively-kernel.org/}} allows short development cycles in combination with direct object
manipulation\cite{Lincke2012Scoping} -- which we dubbed \emph{object-centric development}.

Applications in \ac{lk} can be created by composing basic morphs, which are graphically represented
objects freely placeable and combinable on the screen, and adding functionality to each morph as needed.
Tools like the ObjectEditor and the ObjectInspector allow the developer to manipulate behavior and state of single objects.
Even collaboration works on single objects through the Parts Bin, which is a cloud repository for \ac{lk} objects\cite{Lincke2012LPC}.

In comparison to the relatively abstract class-style development, where the developer implements behavior in an abstract entity: a class,
object-centric development strives to have a concrete connection between the object in development and the tool that manipulates the object.
Through the direct manipulation of objects runtime and development time are merged.
Every change made on an object immediately affects its behavior, which gives the developer the desired early feedback.
This kind of interactive development is particularly good motivated by Bret Victor\cite{victor2012inventing}.

Although it is perfectly possible to develop large tools and applications with object-centric development methods,
unfortunately it is hard to share behavior between different, but similar objects.
Taken a game with multiple characters walking over a map, every character is an object on its own but shares the ability to walk with other characters.
With the current tool support of \ac{lk} the methods that implement walking have to be copied among the characters.
This leads to manually maintained code duplication and therefore to hard to tackle problems with maintainability and code comprehension.



\begin{itemize}
 \item \begin{itemize}
  \item delegation pattern (through morphic scene) is best practice
 \end{itemize}

 \item[II.] However (problem description)
 \begin{itemize}
  \item works in lively for objects that do not share behavior
  \item but its hard to implement shared behavior
  \item pinned down problem is: manually maintained code duplication -- DRY problems
  \item hard to understand large object compositions (e.g. how does the object browser work)
 \end{itemize}

 \item[III.] So what? (problem dimension) 
 \begin{itemize}
   \item typical DRY problems (many changes in code, not maintainable, comprehension -- is this method shared?)
 \end{itemize}

 \item[IV.]  Deficiencies with previous approaches
 \begin{itemize}
  \item delegation: costumes pattern, delegation to one of the owners (problem: opaque objects -- not browsable (comprehension, procedural -- operation on data)
  \item modularization concepts: classes, traits, layer (problem: not concrete, see: parts bin paper)
 \end{itemize}

 \item[IV.] Solution
 \begin{itemize}
  \item use objects, but group them
  \item groups share code (tool-based approach to shared behavior)
  \item make objects visually connected to their tools -- more browsable
 \end{itemize}

\end{itemize}


The remainder of this paper is organized as follows. 
Section \ref{sec:CONCLUSION} concludes this paper.

%%% Local Variables: 
%%% mode: latex
%%% TeX-master: "../paper"
%%% End: 
